\section{Experiments}%#TODO: Needs refactoring 
In this section, we will discuss the basic recipe for "evaluating" in \textbf{software 2.0} or more in general in machine learning.
\\\\
We'll mostly focus on binary classification problems, where the goal is to classify inputs into two categories (for example, "spam" vs "not spam", or "positive" vs "negative").
Usually the metrics used is the proportion of misclassified inputs (also known as \textbf{error}),
meanwhile the proportion of correctly classified inputs is called \textbf{accuracy}.

\subsection{Comparing Models}
Usually, before putting a model into production, we want to compare it with other models to understand which one performs better on a given task.
We could compare different model types but also different parameter settings for the same model type.
\\\\
In simple cases, we can in a 2D dataset plot the predicted output and visually inspect the decision boundary to understand,
but this approach is not scalable to high-dimensional datasets.
Here, the most straightforward approach, is to training both models on the same training set, and pick the one with fewer mistakes.
This generates a few questions:
\begin{itemize}
    \item On which data we compute the error? 
    \item How we eliminate random effects? 
    \item Are the metrics we are using good?
\end{itemize}

\subsection{Test Sets}
As we have seen before, \textbf{never judge your performance on the training set}, since the model has been trained to minimize the error on that specific set, and it may not generalize well to new inputs.
\\\\
So the first thing we do in machine learning is to split the dataset into a training set and a test set.
This should allow us to suppose that the test set is representative of the real-world data distribution.
\\
So in this setting, becames crucial to perform a good split of the dataset.
We don't split the dataset based on proportion, but in the \textbf{absolute size} of the test set, which should be large enough to be representative.
Ideally, we want 10000 examples in the test set, but this is not always possible, especially for small datasets.

\subsubsection{Overfitting in the test set}
However, this approach can lead to overfitting on some cases.
Suppose we want to choose the best hyperparameters for our model, and we are provided with a test set.
We try different hyperparameter settings, and we pick the one that performs best on the test set, supposing that it will generalize well to unseen data.
But, at a certain point, we are provided with a fresh test set, and we find out that the performance is much worse than expected with the chosen hyperparameters.
This is because we have \textbf{overfitted} on the test set.
\\\\
Our method for choosing the best hyperparameters are another learning algorithm, and by testing different hyperparameters on a small test set,
leads us to fit on \textbf{random fluctuations} in the test set, which may not be representative of the real-world data distribution.
\\\\
To simple answer to this problem is \textbf{not to test multiple times on the test set.}

\subsubsection{Validation Sets}
Reusing the test set not only makes you choose the wrong model, but also \textbf{inflates the performance measures}, making them lower than they actually are.
\\\\
To avoid this you need to test which model to use, which hyperparameters to give it, and how to extract features only on the training data.
In order not to do this all on the training set, we can split the training set into a \textbf{training set} and a \textbf{validation set}.
Ideally, the validation set is the same size of your test set, but you can make it a little smaller to get more data for training.
\\\\
Note that this approach by itself does not solve the problem of preventing multiple testing on the test set.
It just provides for a final failsafe to detect it, but you should first convince yourself that the results you see are not down to multiple tests.
You need always to be careful and whenever it's needed to rerun your experiments on training and validation sets.
\\
However, if your dataset is large, and you haven't done anything strange in tuning phase, the results between test set and validation set should be similar.

\subsubsection{Cross Validation}
Doing this split could lead to very small sets, especially if the original dataset is small.
In this case, we can use \textbf{cross validation}, which is a technique that allows us to use the entire dataset for both training and testing, by splitting it into multiple folds.
For example, in a 5-fold cross validation, we split the dataset into 5 equal parts, and we train the model on 4 parts and test it on the remaining part (validation).
We repeat this process 5 times, each time using a different part as the test set,and we average the results to get a more reliable estimate of the model's performance.
\\\\
This can be costly, but ensures that every instance in the dataset is used for training.
After chosen your model and hyperparameters through the cross-validation process, you can test it once on the test set to get an estimate of its performance on unseen data.

\subsubsection{Temporal Data}
If your data has a temporal ordering, you need to take this into account when splitting the dataset.
In this case, you can't sample your test set randomly, since training on samples that are in the future compared to the test set is unrealistic.
You need to maintain the ordering, and split the dataset based on time, using the most recent data as the test set, and the older data as the training set.
\\\\
However, if your data has a timestamp, but there's no meaningful information in the ordering (like in email classification) seeing from the future doesn't usually give you much of an unfair advantage in the task.
\\\\
If you want to apply cross validation to temporal data, that takes a small portion of initial data as the training set and then iteratively adds more data to the training set, while validating on the next portion of data, until you reach the end of the dataset.

\paragraph{Final remark:}
Philosophically talking, when in doubt make sure that the evaluation setting accurately simulates production, and that the validation setting accurately simulates the evaluation setting,

\subsection{Experiment Design}
Now that we now how to experiment, what experiments should we run? Which hyperparameters should we test? 
Basically as long as you are not looking at the test set, you can do whatever to like.

\subsubsection{Hyperparameter Tuning}
When tuning hyperparameters, we have the following options:
\begin{itemize}
    \item \textbf{Trial and Error}: This is the most basic approach, where you simply try different combinations of hyperparameters and see which one performs best on the validation set.
    This is driven by intuition and experience, and it can be time-consuming, especially if you have a large number of hyperparameters to tune.
    \item \textbf{Grid Search}: This is an automated approach, where you define a grid of relevant hyperparameter values, and the algorithm trains and evaluates a model for each combination of hyperparameters in the grid.
    This can be computationally expensive, especially if the grid is large, but it ensures that you are exploring a wide range of hyperparameter combinations.
    \item \textbf{Random Search}: This is another automated approach, where you randomly sample combinations of hyperparameters from a defined distribution, and train and evaluate a model for each combination.
    This can be more efficient than grid search, especially if you have a large number of hyperparameters, since it allows you to explore a wider range of combinations.
\end{itemize}